\documentclass[10pt]{article}
\usepackage[T1]{fontenc}
\usepackage[utf8]{inputenc}
\usepackage[french]{babel}
\usepackage{amsmath,amssymb,mathrsfs,mathpazo}
\usepackage{geometry}
\geometry{a4paper,margin=2cm}
\usepackage{graphicx,array,tabularx,tabularray,colortbl,multirow,multicol,booktabs,enumitem,pifont,ulem,fancybox,diagbox,cancel}
\usepackage{pstricks,pstricks-add,pst-plot,pst-tree,pst-node,pst-text,pst-func,pst-math,pst-bspline,pst-3dplot,pst-eucl,pst-solides3d}
\let\pscalculate\undefined
\usepackage{pst-bezier}
\usepackage{tikz,pgfplots,tkz-tab}
\usetikzlibrary{arrows,arrows.meta,patterns,calc,positioning,trees,decorations.markings,shapes.geometric,shapes.misc,shapes.arrows,intersections,decorations.pathmorphing,decorations.pathreplacing,backgrounds,shadows,babel}
\pgfplotsset{compat=1.18}
\usepgfplotslibrary{fillbetween}
\usepackage[np]{numprint}
\usepackage{listings,makecell,pas-tableur,twoopt}
\let\barre\undefined
\usepackage[pstricks]{tablvar}
\definecolor{darkgreen}{RGB}{0,100,0}
\providecommand{\red}{\color{red}}
\providecommand{\blue}{\color{blue}}
\providecommand{\green}{\color{green}}
\providecommand{\black}{\color{black}}
\providecommand{\white}{\color{white}}
\providecommand{\gray}{\color{gray}}
\providecommand{\cyan}{\color{cyan}}
\providecommand{\magenta}{\color{magenta}}
\providecommand{\yellow}{\color{yellow}}
\providecommand{\violet}{\color{violet}}
\providecommand{\orange}{\color{orange}}
\providecommand{\euro}{\texteuro}
\providecommand{\EUR}{\texteuro}
\providecommand{\ieme}{\textsuperscript{e}}
\providecommand{\degres}{\ensuremath{^\circ}}
\providecommand{\degree}{\ensuremath{^\circ}}
\providecommand{\vv}[1]{\overrightarrow{#1}}
\providecommand{\Coord}[1]{#1}
\providecommand{\decosix}{}
\providecommand{\starredbullet}{\ensuremath{\bullet}}
\pagestyle{empty}
\setlength{\parindent}{0pt}
\usepackage{adjustbox}\DeclareUnicodeCharacter{2236}{:}\DeclareUnicodeCharacter{2190}{\ensuremath{\leftarrow}}
\providecommand{\R}{\mathbb{R}}
\providecommand{\N}{\mathbb{N}}
\providecommand{\Z}{\mathbb{Z}}
\providecommand{\Q}{\mathbb{Q}}
\providecommand{\D}{\mathbb{D}}
\providecommand{\C}{\mathbb{C}}
\providecommand{\e}{{\mathrm{e}}}
\providecommand{\vect}[1]{\overrightarrow{#1}}
\providecommand{\vectt}[1]{\overrightarrow{\mathrm{#1}}}
\providecommand{\vvt}[1]{\overrightarrow{\mathrm{#1}}}
\providecommand{\np}[2][]{{#2~#1}}
\providecommand{\vv}[1]{{\overrightarrow{#1}}}
\renewcommand{\d}{{\mathrm{d}}}
\providecommand{\degres}{{^\circ}}
\providecommand{\degre}{{^\circ}}
\providecommand{\blue}{\color{blue}}
\providecommand{\red}{\color{red}}
\providecommand{\enskip}{\quad}
\providecommand{\llbracket}{\lbrack\!\lbrack}
\providecommand{\rrbracket}{\rbrack\!\rbrack}
\providecommand{\euro}{{\text{€}}}
\providecommand{\no}{{\text{n°}}}
\providecommand{\interval}[2]{{[#1;#2]}}
\providecommand{\intervFO}[2]{{[#1;#2[}}
\providecommand{\dotfill}{\ldots}

\providecommand{\ds}{\displaystyle}
\providecommand{\textding}[1]{#1}
\providecommand{\Oij}{\(\mathrm{O};\vec{\imath},\vec{\jmath}\)}
\providecommand{\Oijk}{\(\mathrm{O};\vec{\imath},\vec{\jmath},\vec{k}\)}
\providecommand{\Ouv}{\(\mathrm{O};\vec{u},\vec{v}\)}
\providecommand{\intervFF}[2]{\ensuremath{[#1;#2]}}
\renewcommand{\euro}{\ifmmode\text{\texteuro}\else\texteuro\fi}
\DeclareUnicodeCharacter{00B0}{\ifmmode^{\circ}\else\textdegree\fi}
\DeclareUnicodeCharacter{20AC}{\ifmmode\text{\texteuro}\else\texteuro\fi}
% Énoncés : Éducation nationale. Transcription et corrigés : APMEP, Denis Vergès et contributeurs. Découpage et classement : CoopMaths, Sébastien Lozano, Éric Elter et contributeurs. Adaptation de la présentation : Maths Explorer.
% CC BY-SA : https://coopmaths.fr/www/about/
% https://coopmaths.fr/alea/static/bac/2025/tex/bac_2025_06_sujet2_asie_2.tex
\begin{document}
\noindent\textit{APMEP / CoopMaths. CC BY-SA. Présentation adaptée par Maths Explorer.}\par\medskip
\medskip

\textbf{Partie A}

\medskip

Soit $\left(u_{n}\right)$ la suite définie par $u_{0}=30$ et, pour tout entier naturel $n, u_{n+1}=\dfrac{1}{2} u_{n}+10$.\\
Soit $\left(v_{n}\right)$ la suite définie pour tout entier naturel $n$ par $v_{n}=u_{n}-20$.

\begin{enumerate}
\item Calculer les valeurs exactes de $u_{1}$ et $u_{2}$.
\item Démontrer que la suite $\left(v_{n}\right)$ est géométrique de raison $\dfrac12$,
\item Exprimer $v_{n}$ en fonction de $n$ pour tout $n$ entier naturel.
\item En déduire que, pour tout entier naturel $n,\,  u_{n}=20+10\left(\dfrac{1}{2}\right)^{n}$.
\item Déterminer la limite de la suite $\left(u_{n}\right)$. Justifier la réponse.
\end{enumerate}

\textbf{Partie B}

\medskip

Soit $\left(w_{n}\right)$ la suite définie pour tout entier naturel $n$ par :

\[
\begin{cases}
w_{0}=45 \\
w_{n+1}=\dfrac{1}{2} w_{n}+\dfrac{1}{2} u_{n}+7
\end{cases}
\]

\begin{enumerate}
\item Montrer que $w_{1} = 44,5$.
\end{enumerate}

On souhaite écrire une fonction suite, en langage Python, qui renvoie la valeur du terme $w_{n}$ pour une valeur de $n$ donnée. On donne ci-dessous une proposition pour cette fonction suite.

\begin{center}

\begin{tabular}{c| l|}\cline{2-2}
1&def suite(n):\\
2&\quad U=30\\
3&\quad W=45\\
4&\quad for i in range (1,n+1):\\
5&\quad\quad   U=U/2+10\\
6&\quad\quad  W=W/2+U/2+7\\
7&\quad return W\\ \cline{2-2}

\end{tabular}
\end{center}

\begin{enumerate}[resume]

\item L'exécution de suite(1) ne renvoie pas le terme $w_{1}$. Comment modifier la fonction suite afin que l'exécution de suite(n) renvoie la valeur du terme $w_{n}$ ?
\item 
\begin{enumerate}
	\item Montrer, par récurrence sur $n$, que pour tout entier naturel $n$ on a :

\[w_{n}=10 n\left(\dfrac{1}{2}\right)^{n}+11\left(\dfrac{1}{2}\right)^{n}+ 34\]

\item  On admet que pour tout entier naturel $n \geqslant 4$, on a : $0 \leqslant 10 n\left(\dfrac{1}{2}\right)^{n} \leqslant \dfrac{10}{n}$.

Que peut-on en déduire quant à la convergence de la suite $\left(w_{n}\right)$ ?
\end{enumerate}
\end{enumerate}

\bigskip

\end{document}
